\documentclass[10pt,a4paper]{amsart}
\usepackage[margin=19mm]{geometry}
\usepackage{amsmath,amssymb,mathtools}
\usepackage[scr=rsfs,cal=euler]{mathalfa}
\usepackage{xfrac}
\usepackage[unicode,hidelinks,bookmarks=false]{hyperref}
\newcommand{\ie}{i.e.\ }
\newcommand{\cf}{cf.\ }
\newcommand{\resp}{respectively\ }
\newcommand{\Subsection}{\subsection}
\providecommand{\nobreakdash}{\nobreak\hbox{-}}
\setlength{\emergencystretch}{2em}
\multlinegap0pt
\usepackage{bm}
\usepackage{bbm}
\usepackage{dsfont}
\usepackage{xspace}
\usepackage{cancel}
\usepackage{mathtools}
\usepackage{cleveref}
\usepackage{amsthm}
\makeatletter
\@ifundefined{claim}{\newtheorem{claim}{Claim}}{}
\@ifundefined{claimproof}{\newtheorem{claimproof}{Claimproof}}{}
\@ifundefined{theorem}{\newtheorem{theorem}{Theorem}}{}
\@ifundefined{lemma}{\newtheorem{lemma}[theorem]{Lemma}}{}
\makeatother
\makeatletter
\expandafter\ifx\csname claimproof\endcsname\relax
\newenvironment{claimproof}[1][Proof]{
  \renewcommand{\oldqed}{\qedsymbol}
  \renewcommand{\qedsymbol}{\endofClaim}
  \begin{proof}[#1]}
\fi
\makeatother
\title[Range of Clique Counts in Graphs]{Range of Clique Counts in Graphs}
\author{Mihir Neve, Alexey Pokrovskiy, Andrey Shapiro}
\date{Source: arXiv:2609.21739v1 (CC BY 4.0)}
\begin{document}
\maketitle

\noindent\emph{Source and licence.} Range of Clique Counts in Graphs. Version arXiv:2609.21739v1 (CC BY 4.0), distributed under the Creative Commons Attribution 4.0 International licence, as recorded in the arXiv licence metadata for this exact version and shown on the abstract page https://arxiv.org/abs/2609.21739v1. Full rights notice: licenses/arxiv-2609.21739-CC-BY-4.0.txt.\par\smallskip

\noindent\textit{Editorial scope.} Counted unit: the Theorem whose source label is \texttt{None} in the article (source0.tex lines 349--351, source label \texttt{None}), delivered together with the article's own complete original proof of it (source0.tex lines 353--438, one \texttt{proof} environment of 86 lines). No mathematics is written, repaired or re-derived: statement and proof are the article's own text line for line. Required context retained uncounted: thm:main (lines 118--121); q2 (lines 135--144). Every definition, display and label of those lines is retained unchanged. Omitted from this package and not redistributed: 1--117, 122--134, 145--348, 352--352, 439--453. Nothing inside a retained excerpt is omitted or altered: every retained body is delivered exactly as the article's lines are written, so no omission receipt of any kind accompanies this package. Citations: the retained excerpts carry no citation command at all, so no bibliography entry of the article is transcribed and no reference is redistributed. Reference adaptation: none was needed and none was made; every reference inside the delivered unit resolves inside it, so no unresolved reference or citation is produced. Printed slips kept literal: no printed word, symbol, hypothesis or number of the counted statement or of its proof was altered. Layout and class: the article's own class is \texttt{article} with packages including geometry, amsmath, amssymb, amsthm, amsbsy, mathtools, bbm, dsfont, enumitem, microtype, parskip, graphicx, xcolor, mathrsfs; the wrapper is the lane's pinned \texttt{amsart} editorial wrapper at 10pt on A4, which restates the theorem-like environments the retained text needs and adds the article's own macro definitions that the retained text needs (none), with extra wrapper packages (bm, bbm, dsfont, xspace, cancel, mathtools, cleveref). The delivered numbering is the unit's own. Assets: no retained span contains a figure environment, a graphics-inclusion command, a PDF-page inclusion, a table or a TikZ picture; no image file is copied into this package, the package declares no asset dependency and no image file is redistributed with it. Licence: version arXiv:2609.21739v1 of this article is distributed under the Creative Commons Attribution 4.0 International licence, as recorded in the arXiv licence metadata for this exact version and shown on the abstract page https://arxiv.org/abs/2609.21739v1. A full rights notice is delivered at licenses/arxiv-2609.21739-CC-BY-4.0.txt. Characters: every retained excerpt is pure ASCII, so the delivered engine, XeLaTeX with its default Latin Modern text font, reports no missing character.
\begin{theorem}
\label{thm:main}
    For sufficiently large $n$, we have $|\Gamma(n)| \geq 2^{n-4\ln(2)\log^3(n)}$.
\end{theorem}

\begin{lemma}\label{q2}

For every non-negative integer $a$, and positive integers $\ell$ and $m$ satisfying $a + (\ell+3)2^\ell\leq m$, there exist a constant~ $C = C_{a,\ell,m}$ and a function~$f = f_{\ell,m}\colon \bigl[2^{2^\ell}\bigr]\to \mathbb{Z}_{\ge 0}$ such that for all $x
\in \bigl[2^{2^\ell}\bigr]$, there exists a graph $H_x$ on $m + 2\ell$ vertices satisfying the following properties.

\begin{enumerate}
    \item There exists a subset $M_0\subseteq V(H_x)$ of size $m$ that forms a clique, and 
    \item we have $\gamma(H_x)=C + 2^a (x-1) + 2^{a+2^\ell+1}f(x)$. \label{clique_number}
\end{enumerate}
\end{lemma}

\begin{theorem}
    For all sufficiently large $n$, define $m = n-\lfloor 4\ln(2)\log^3n\rfloor +8$. Then for each integer~$x \in \bigl[2^{m-8}\bigr]$, there exists an $n$-vertex graph~$G_x$ such that for all distinct~$x, y\in \bigl[2^{m-8}\bigr]$, we have~$\gamma(G_x)\neq \gamma(G_y)$. 
\end{theorem}

\begin{proof}
    \setcounter{equation}{0}
    Let $n, m$ be given as in the statement of the theorem. We begin by constructing the intervals $[a_1,a_2), [a_2,a_3),\dots,[a_w,a_{w+1})$ as follows. Set $a_1\coloneqq0$. Having defined $a_i\le m-8$, let $\ell_i$ be the largest positive integer satisfying $a_{i}+(\ell_i+3)2^{\ell_i}\leq m$, and set     
    $a_{i+1} \coloneqq a_{i}+2^{\ell_i}$. Stop at the first $w$ for  which $a_{w+1}>m-8$. Since $a_i\le m-8$, the choice $\ell=1$ is admissible, so $\ell_i$ is well-defined and satisfies $\ell_i \geq 1$. Also, trivially  $\ell_i\leq \log m$. We shall use~$\Lambda_i\coloneqq a_{i+1}-a_i = 2^{\ell_i}$ to denote the length of the interval $[a_i, a_{i+1})$ for each $i \in [w]$.

    We construct $G_x$ for every $x\in[2^{a_{w+1}}]$.
Since $a_{w+1}>m-8$, restricting to
$x\in[2^{m-8}]$ proves the theorem. Moreover, each integer $x \in [2^{a_{w+1}}]$ can be uniquely associated with a $w$-tuple $(x_1, x_2, \dots, x_w)$ such that 
    \begin{equation}
        x-1=\sum_{i \in [w]} 2^{a_i}(x_i-1)\quad \text{and} \quad x_i\in \bigl[2^{a_{i+1}-a_i}\bigr] = \bigl[2^{\Lambda_i}\bigr] \text{ for each $i \in [w]$.}\label{eqTH1}
    \end{equation} 
    
    We construct the required $n$-vertex graph~$G_x$ as follows. Let $(x_1, x_2, \dots, x_w)$ be the unique $w$-tuple associated with~$x$, as described above. Note that by construction of the intervals~$[a_i, a_{i+1})$, we have that $a_i + (\ell_i + 3)2^{\ell_i} \leq m$ for each~$i \in [w]$. Hence, for each~$i \in [w]$, \Cref{q2} (with parameters~$a_i, \ell_i,m$) yields a constant~$C_i$ and a function~$f_i:\bigl[2^{\Lambda_i}\bigr]\to \mathbb{Z}_{\ge0}$ such that for the given~$x_i\in \bigl[2^{\Lambda_i}\bigr]$, there exists an $(m+2\ell_i)$-vertex graph~$H_{x_i}$ with $ \gamma(H_{x_i}) = C_i+2^{a_i}(x_i-1)+2^{a_i+\Lambda_i+1}f_i(x_i)$.

    Recall that for each $i \in [w]$, the graph $H_{x_i}$ contains a subset $M^i_0\subseteq V(H_{x_i})$ of size $m$ that forms a clique in $H_{x_i}$. We let $H_x$ be the clique sum of the graphs $H_{x_1}, \dots, H_{x_w}$ with respect to the cliques $M_0^1, \dots, M_0^w$. Observe that the graph~$H_x$ has exactly  $N \coloneqq m+2\sum_{i=1}^w \ell_i$ vertices, and the quantity~$N$ is independent of the choice of~$x$. We shall prove in the latter half of this proof that $N \leq n$. Assuming this inequality, the required $n$-vertex graph~$G_x$ is obtained from $H_x$ by additionally adding $n-N$ isolated vertices to $H_x$. 

    Next, we compute the number of cliques~$\gamma(G_x)$ in the graph~$G_x$. Observe that forming this clique sum of the graphs~$H_{x_1}, \dots, H_{x_w}$ does not create any new cliques. However, since the clique sum is constructed by superimposing the cliques~$M_0^1, \dots, M_0^w$ onto a common clique~$M_0 \subseteq V(H_x)$ of size~$m$, each of the~$2^m$ cliques within~$M_0$ is overcounted precisely $w-1$ times. So, it follows that $\gamma(H_x) = \sum_{i =1}^w \gamma(H_{x_i}) - (w-1)2^m$. Including the $n -N$ cliques induced by the set of isolated vertices~$V(G_x) \setminus  V(H_x)$, and by equation~\eqref{eqTH1}, we have that
    $$\gamma(G_x) = \sum_{i\leq w} \gamma(H_{x_i})  - (w-1)2^m + n-N = C+x+\sum_{i\leq w}2^{1+a_{i+1}}f_i(x_i),$$
    where we define $C\coloneqq \sum_{i\leq w} C_i - (w-1)2^m + n - N - 1$. We now show that the mapping~$x \mapsto \gamma(G_x)$ on the domain~$[2^{a_{w+1}}]$ is an injection.
    
    \begin{claim}\label{clm:injection}
        $\gamma(G_x)\neq \gamma(G_y)$ for all distinct $x,y\in [2^{a_{w+1}}]$.
    \end{claim}
    
    \begin{claimproof}
        Let distinct $x, y \in [2^{a_{w+1}}]$ be given. Let $(x_1, \dots, x_w)$ and $(y_1, \dots, y_w)$ be the unique $w$-tuples satisfying equation~\eqref{eqTH1}  for~$x$ and~$y$, respectively. Finally, let~$k \in [w]$ be the smallest integer for which we have~$x_k \neq y_k$. Note that such an integer~$k$ exists by the assumption~$x \neq y$. Now, evaluating $\gamma(G_z)$ modulo $2^{a_{k+1}}$ for $z \in \{x, y\}$, we obtain 
        \begin{equation}
            \gamma (G_z) \bmod{2^{a_{k+1}}} = \Biggl[\, C + 1 + \sum_{i = 1}^k 2^{a_i}(z_i -1) + \sum_{i = 1}^{k-1} 2^{1 + a_{i+1}} f_i(z_i) \Biggr] \bmod{2^{a_{k+1}}}. \label{eqTH2}
        \end{equation}

        By the minimality of $k$, we have that $x_i = y_i$ for all $i \in [k-1]$; and consequently, it follows that~$f_i(x_i) = f_i(y_i)$ for all $i \in [k-1]$. Thus, by equation~\eqref{eqTH2}, we have 
        \begin{align*}
            \bigl(\gamma(G_x) - \gamma(G_y) \bigr) \bmod 2^{a_{k+1}}&=\Biggl[\,\sum_{i = 1}^k 2^{a_i}(x_i -y_i) + \sum_{i = 1}^{k-1} 2^{1 + a_{i+1}} \bigl(f_i(x_i) - f_i(y_i)\bigr) \Biggr] \bmod 2^{a_{k+1}}\\[0.5em]
            &= 2^{a_k} (x_k - y_k) \bmod 2^{a_{k+1}}\, \neq\, 0.
        \end{align*}
        
        Here, the last relation follows from the fact that~$x_k \neq y_k$ and that $\bigl|2^{a_k}(x_k -y_k)\bigr| < 2^{a_{k+1}}$ since we have~$x_k , y_k \in \bigl[2^{a_{k+1} - a_k}\bigr]$. Thus, as $\gamma(G_x) \bmod{2^{a_{k+1}}} \neq \gamma(G_y) \bmod{2^{a_{k+1}}}$, it follows that~$\gamma(G_x) \neq \gamma(G_y)$, concluding the proof of the claim. 
    \end{claimproof}
     
     As $a_{w+1} > m - 8$ by definition, Claim~\ref{clm:injection} implies that $|\Gamma(n)| \geq 2^{m-8}=2^{n-\lfloor4\ln(2)\log^3 n\rfloor}\ge 2^{n-4\ln(2)\log^3 n}$, as witnessed by the collection of $n$-vertex graphs $\bigl\{G_x : x \in [2^{a_{w+1}}] \bigr\}$. However, it remains to show that the number of vertices in the graph $H_x$ for $x \in [2^{a_{w+1}}]$, as given by $N = m + 2\sum_{i = 1}^w \ell_i$, satisfies $N \leq n$. For this, we observe that the inequality $\ell_i \leq \log m$ for each $i \in [w]$ implies that $N \leq m + 2w \log{m}$, and hence, it suffices to prove that $m + 2w\log m \leq n$ for sufficiently large $n$. 
     
     We begin by obtaining an upper bound on the number of intervals~$w$. Let us do this by considering a different sequence~$b_1, \dots, b_{w'}$, defined using a non-integer relaxation of the recurrence relation used to define the sequence~$a_1, \dots, a_w$. More precisely, set~$b_1 \coloneqq 0$. For each~$i \geq 1$, let~$\ell_i'$ be the (not necessarily integer) solution to the equation $b_i+(2\log m)2^{\ell_i'}=m$, and set~$b_{i+1} \coloneqq b_i+2^{\ell'_i}$. We let~$w'$ be the smallest integer such that~$b_{w'+1}> m-8$. We first prove that~$w\leq w'$, and then we obtain an upper bound on~$w'$.

     \begin{claim}
     \label{clm:wwprime_reln}
     For sufficiently large~$n$, we have that $w \leq w'$.
     \end{claim}

     \begin{claimproof}

     Let $n$ be large enough so that $\log m > 17$. We first prove that $a_k \geq b_k$ for each $k \in [w+1]$. This is done by induction, starting with the base case $a_1 = b_1 = 0$. Suppose that $a_i\geq b_i$ for some $i\leq w$. If we have $2^{\ell_i} \geq 2^{\ell_i'}$, then the inductive hypothesis implies that $$a_{i+1}=a_i +2^{\ell_i}\geq b_i+2^{\ell_i'}=b_{i+1},$$
     as required. Hence, we may assume that $2^{\ell_i} < 2^{\ell_i'}$. 
     
     Recall that $\ell_i$ is the largest positive integer for which we have $a_i + (\ell_i + 3)2^{\ell_i} \leq m$. It is easy to see that $\ell_i\leq \log m-4$, as otherwise we obtain the inequality  
     $$ a_i+(\ell_i+3)2^{\ell_i}\geq  (\log (m)-1) \frac m {16} > m,$$
     thereby contradicting the definition of $\ell_i$. 
     
     The maximality of $\ell_i$ in the above definition also implies that $a_i + (\ell_i + 4)2^{\ell_i+1} > m$. Combining this fact with the recurrence relation for $a_{i+1}$ yields the following inequality.
     \begin{equation}
         a_{i+1} = a_i + 2^{\ell_i} > m - (2\ell_i + 8)2^{\ell_i} + 2^{\ell_i} = m - (2\ell_i + 7)2^{\ell_i}.
         \label{eqTH3}
     \end{equation}
     Similarly, the definition of $\ell_i'$ and the recurrence relation for $b_{i+1}$ results in the relation
     \begin{equation}
         b_{i+1} = b_i + 2^{\ell_i'} = m - (2\log m)2^{\ell_i'} + 2^{\ell_i'}= m - (2\log m-1)2^{\ell_i'}.
         \label{eqTH4}
     \end{equation}
     
     Now, the inequality $a_{i+1} \geq b_{i+1}$ follows by combining equations \eqref{eqTH3} and \eqref{eqTH4} to obtain $$a_{i+1}-b_{i+1}>(2\log m-1)2^{\ell_i'}-(2\ell_i+7)2^{\ell_i}>(2\log m-2\ell_i-8)2^{\ell_i'} \geq 0,$$
     where, for the second and third inequalities, we make use of the assumption $2^{\ell_i} < 2^{\ell_i'}$ and the fact $\ell_i\leq \log m-4$, respectively. 
     
     Thus, we have $a_k \geq b_k$ for all $k \in [w+1]$. The claim now follows. Indeed, if $w \geq w' + 1$, then we must have $a_w \geq a_{w' +1} \geq b_{w'+1} > m-8$, thereby contradicting the definition of $w$.
     \end{claimproof}
    
     Let us now obtain an upper bound on~$w'$. Set~$q \coloneqq 2\log m$. Then, by the definition of~$\ell_i'$, we have that~$2^{\ell_i'} = (m - b_i)/q$. Substituting this in the equation~$b_{i+1} = b_i + 2^{\ell_i'}$, we obtain the recurrence relation $b_{i+1}=\mu b_i + m/q$, where $\mu \coloneqq (q-1)/q$.
     Note that any sequence $\{s_i\}_{i \in \mathbb{N}}$ satisfying the recurrence relation $s_{i+1} = \alpha s_i + \beta$ with $\alpha \neq 1$ and $s_1 = 0$ has a closed form given by $s_{k+1} = \beta (1 - \alpha^k)/(1 -\alpha)$. Hence, it follows that 
     $$b_{w'}=\frac{m}{q}\cdot \frac{1-{\mu}^{w'-1}}{1-\mu} = m \bigl(1 - \mu^{w' - 1}\bigr).$$ 
     
     As $b_{w'} \leq m - 8$ by definition of $w'$, we have that $m \mu^{w' -1} \geq 8$, or equivalently, $(1/\mu)^{w' - 1} \leq m/8$. By taking natural logarithm and solving for $w'$, we obtain the inequality 
     \begin{eqnarray*}
     w'-1 &\leq& \frac{\ln(m/8)}{\ln(1/\mu)} = \frac{\ln (m/8)}{\ln\bigl(1 + 1/(q-1)\bigr)} \leq \frac{1+1/(q-1)}{1/(q-1)}\cdot \ln( m/8)\\
     &=& q \cdot \ln(2) \log(m/8)= 2\ln(2)\log^2(m)-2\ln(2)\log(8)\log(m)\\[0.1em]
     &\leq& 2 \ln(2) \log^2 (m) - 2,    
     \end{eqnarray*}
    where, the second inequality uses the fact that $\ln(1+x) \geq x/(1+x)$ for all $x > -1$. Thus, along with Claim~\ref{clm:wwprime_reln}, we obtain the upper bound $w\leq 2\ln(2)\log^2(m)  -1 $.
    The proof of Theorem~\ref{thm:main} now follows by observing that for sufficiently large~$n$, we have $$N \leq m+2w\log(m) \leq  n-4\ln(2)\log^3(n) +9 + 4\ln(2)\log^3(m)-2\log(m) \leq n.$$\end{proof}

\end{document}
